\chapter{``\ldots{} Todo o tempo do mundo''}

\lettrine{E}{nquanto nos sentamos} aqui durante este retiro, temos que prestar atenção a
coisas que não são de todo interessantes. Podem até ser desagradáveis e
dolorosas. Perseverar pacientemente ao invés de sairmos a correr em busca de
algo interessante é uma boa disciplina, não é? É bom sermos capazes de
simplesmente aguentar o tédio, a dor, a raiva, a ganância, todas essas coisas,
em vez de fugirmos sempre delas\ldots{} Paciência é uma virtude tão importante.
Se não tivermos paciência, não haverá possibilidade alguma de nos iluminarmos.
Sejamos \emph{extremamente} pacientes\ldots{}

Eu gostava do tipo de meditação em que me sentava e ficava muito calmo -- e,
então, quando a dor surgia no corpo, o que eu queria era livrar"-me dela para
poder ficar naquele estado de calma. Então comecei a ver que querer
\emph{livrar"-me} da dor, era um estado de espírito miserável. Às vezes, ficamos
sentados por várias horas; outras vezes a noite inteira. Podemos fugir disso,
mas passado um tempo começamos a saber lidar com a dor física. Cheguei a usar
práticas como ``ter todo o tempo do mundo para estar com a dor'', em vez de dar
o meu melhor para me livrar dela de forma a poder voltar à minha ``verdadeira''
meditação. Se as dores no meu corpo surgirem na consciência, eu tiro um tempo
para lidar com elas, em vez de pensar, ``Como é que eu me posso livrar delas
para obter alguma felicidade?''.

De certa forma, quando eu dizia, ``Tenho todo o tempo do mundo, o resto da minha
vida para estar com esta dor'', isso parava a tendência de me querer livrar
dela. A minha mente ficava na realidade mais lenta por longos períodos de tempo,
sem seguir ou criar um desejo. Alguns de vós têm essa ideia de \emph{vencer} a
dor, de ultrapassar o ``limiar da dor'' -- mas isso é um desastre\ldots{}
